\section{Technical construction}

\subsection{A mandatory causal state}

% claims: C02 C03; sources: S03 S04
For $x=(a,b,c,d)\in\{0,\ldots,15\}^4$, the task is
\begin{equation}
z_1=(a+b)\bmod16,\quad \ztwo=(z_1+c)\bmod16,\quad
\zthree=(5\ztwo+d)\bmod16,\quad y=\zthree\bmod8.
\end{equation}
Operationally, the model must predict a 16-way hard state \ztwo{}, combine only
that state with $d$ to predict hard \zthree{}, and use only \zthree{} to predict
the eight-way answer. The downstream modules receive no true labels, soft
probabilities, logits, residual input features, prompt features, or shared
cache. We call \ztwo{} a \emph{causal state} because registered interventions on
it change the learned downstream computation, not merely because it is
decodable.

% claims: C03; sources: S03 S04
The alias intervention sets $\ztwo'=\ztwo+8\bmod16$. Since
$5\cdot8=40\equiv8\pmod{16}$, it also shifts \zthree{} by eight; because
$y=\zthree\bmod8$, the final answer is unchanged. By contrast,
$\ztwo'=\ztwo+1$ changes the registered answer. After either intervention, the
learned transition still produces \zthree{} rather than receiving it from the
controller. Same-state, sham, wrong-state, and direct-\zthree{} controls test
additional failure modes.

% claims: C02 C03 C12; sources: S03 S04 S18
The study prospectively selected Model 1 at the first checkpoint satisfying its
frozen gates. The model has 12,648 parameters and supervised 16-way hard
states. A separate causal confirmation tested 256 matched units, producing 512
oriented examples across seven registered intervention families. These tests
establish correct transport for the selected model only. The wider
three-seed recipe failed its preregistered three-of-three requirement: Models 1
and 3 passed, while Model 2 failed a frozen subgroup gate.

\subsection{Authenticated evidence}

% claims: C04 C05; sources: S05 S06 S13
Trusted instrumentation issues \hmac-SHA256 receipts after observing the hard
state boundaries. An exact receipt binds one consumed state event to its run,
context, execution, sequence, and intervention. An abstract receipt links that
event to a reusable category in a frozen ontology and model/checkpoint domain.
A chain receipt checks the registered order and relations
$\ztwo\rightarrow\zthree\rightarrow y$. Generation is authorised only after
the required exact, abstract, and chain evidence passes.

\begin{table}[t]
% claims: C04 C05; sources: S05 S06 S13
\caption{Exact-event evidence is not abstract-state or carrier evidence.}
\label{tab:provenance}
\centering
\small
\begin{tabularx}{\linewidth}{@{}Yccc@{}}
\toprule
Candidate evidence & Exact & Abstract & Carrier claim \\
\midrule
Correct execution and state & \Accept & \Accept & compatible \\
Different execution, same state & \Reject & \Accept & compatible \\
Different registered state & \Reject & \Reject & incompatible \\
Required receipt missing & \Abstain & \Abstain & none \\
Present but invalid receipt & \Reject & \Reject & none \\
\bottomrule
\end{tabularx}
\end{table}


% claims: C04 C05; sources: S05 S06
The trusted computing base contains the checkpoint, state-boundary
instrumentation, receipt issuer, verifier/controller, and key custody. The
registered tests cover missing, malformed, crossed, replayed, reordered, and
wrong-state evidence across contexts and executions. They do not cover a
malicious issuer holding the master key, an arbitrary controller, key
compromise, adaptive rewriting, or an untrusted verifier.

\subsection{From state class to statistical carrier}

% claims: C04 C05; sources: S05 S06 S12 S13
After verification, the controller derives a \emph{carrier identity}---called
an authority in the implementation---from the ontology, model/checkpoint
domain, state type, and abstract category. It deliberately excludes receipt
bytes and exact execution identity. Consequently, two executions in the same
state class share an authority, while distinct state classes select different
authorities.

% claims: C01; sources: S01 S10
The surface realiser emits a fixed four-field object through a finite grammar.
Of 49 autoregressive steps, 24 admit eight allowed alternatives. At an eligible
position, the state authority defines a keyed ``green'' set of four
alternatives, and the realiser adds a fixed logit bias of 2 to those
alternatives before sampling. This SynthID-style mechanism refers only to
statistical in-generation carriage; it is not equivalent to a deployed
SynthID system.

% claims: C06; sources: S07
To isolate the matched comparison, the sampler uses
\emph{position-indexed randomness}. Each output position receives a
deterministic uniform draw indexed by the pair, output, step, and subdraw, but
not by the state, answer, receipt, or generated history. Token probabilities
remain autoregressive and therefore history-dependent. The indexed uniforms
nevertheless let natural and alias executions reuse the same primitive draws
after their token histories diverge. The implementation name for this sampler
is \rngtwo{}.

% claims: C07; sources: S08 S09
For output $i$, let $G_i(A)$ be the number of the 24 eligible positions that
fall in authority $A$'s green set. The per-output score and eight-output
aggregate are
\begin{equation}
Z_i(A)=\frac{G_i(A)-12}{\sqrt{6}},\qquad
S_A=\frac{1}{\sqrt{8}}\sum_{i=1}^{8}Z_i(A).
\end{equation}
The first expression measures excess green choices relative to the null
expectation; the second combines eight generated outputs for one execution.
The intended authority is accepted only if its aggregate exceeds the frozen
threshold $T=4.04145$, uniquely ranks first among 16 candidates, and has a
strictly positive margin over the strongest alternative. Calibration fixed
these rules before validation.

% claims: C09; sources: S10 S12 S13
Figure~\ref{fig:contrast} is the integrated test of this construction. It does
not claim that literally every field is fixed: authenticated \ztwo{},
model-generated \zthree{}, receipt context, authority-conditioned
probabilities, and the intended signal change downstream. It does control the
prompt, final answer, printed payload, configuration, output indices, and
primitive random draws within the frozen pipeline.
